\section{The Vegas Language}
\label{sec:language}

This section introduces the surface language by example, fixes
intuitive semantics for each construct, and gives the small grammar
in BNF.  The next section traces what happens to a program once the
compiler takes it apart.

\subsection{Programs and roles}

A Vegas program declares optional type aliases and macros, then a
single top-level \TT{game} block.  The block names the roles, opens a
\TT{join} phase that locks each role's deposit into a common pool,
and proceeds with a sequence of action statements until a
\TT{withdraw} clause splits the pool.

\begin{lstlisting}[style=vg]
type bid = {0 .. 10}

game main() {
  join Seller() $ 0
       Bidder() $ 100;
  // ... statements ...
  withdraw { Seller -> ...; Bidder -> ... }
}
\end{lstlisting}

Role identifiers start with an uppercase letter; field names with
lowercase.  A field is always identified relative to its owning
role: \TT{Bidder.b} denotes the field \TT{b} of role \TT{Bidder}.
Each field has at most one writing action; the language has no
explicit reassignment.

\subsection{Action statements}

There are five action kinds.  Three of them (\TT{yield},
\TT{commit}, \TT{reveal}) write one or more fields owned by
the actor; \TT{join} and \TT{random} declare a role
(strategic and chance, respectively) and record its deposit,
without writing protocol fields.

\paragraph{\TT{join}.}  Initial deposit and admission.  The
\TT{join} phase appears once, at the top of the game.  A role listed in the join phase is bound to the game, with its
deposit locked into the pool.

\paragraph{\TT{yield}.}  A public move.  The actor declares a value
of a given type:
\begin{lstlisting}[style=vg]
yield Guest(d: door);
\end{lstlisting}
After this action fires, \TT{Guest.d} is visible to every role.  A
\TT{yield} can carry a \TT{where} guard, e.g.
\TT{yield Host(goat: door) where Host.goat != Guest.d}; this guard
is enforced when the move is played.

\paragraph{\TT{commit}.}  A private move.  The actor declares a
value that is held cryptographically (in the deployed contract:
under a hash with salt) until a later \TT{reveal}:
\begin{lstlisting}[style=vg]
commit Host(car: door);
// ...
reveal Host(car: door) where Host.goat != Host.car;
\end{lstlisting}
A \TT{where} clause on a \TT{commit} is \emph{deferred}: it is
checked at the matching \TT{reveal}, not at the commit, because the
committed value is not yet observable.

\paragraph{\TT{reveal}.}  Make a previously committed field
public.  Every \TT{commit} is paired with exactly one \TT{reveal}
on each non-quit path through the game; the type checker enforces
this (a committed value that is never revealed and never reachable
through a quit handler is rejected).  Before reveal, only the
committing actor knows the value --- operationally, the contract
holds only the hash.

\paragraph{Randomness.}  Vegas distinguishes two surface forms,
matching the two operational modes available on EVM today.

A \TT{random} declaration introduces a chance role with an actor
identity but no deposit: useful when the draw is \emph{private}
and must be committed and later revealed by a designated actor.
\begin{lstlisting}[style=vg]
random Host;
commit Host(car: door ~ uniform { 0, 1, 2 });
reveal Host(car: door);
\end{lstlisting}
A \TT{sample} declaration introduces an \emph{anonymous public}
draw with no actor:
\begin{lstlisting}[style=vg]
sample (w: ticket);
\end{lstlisting}
Anonymous samples live under a reserved label \TT{Sample}; their
value is computed on chain from \TT{block.prevrandao} with
domain separation.  No actor identity is required; the field is
referenced in expressions as \TT{Sample.w}.

A parameter may carry a \emph{distribution annotation}
\TT{$\sim$ \textit{D}} declaring its analysis prior.  Two surface
forms are available, both finite-support and rational-weighted:
\TT{uniform \{ \dots \}} and
\TT{weighted \{ v\textsubscript{1}: w\textsubscript{1}, \dots \}}.
The annotation is an \emph{analysis assumption}: the analysis
backends (Gambit, MAID) interpret the action as a draw from the
declared distribution.  On a strategic role's action it asserts
that the role plays through; on a chance role's or sample's
action it overrides the default uniform-over-domain prior.

\paragraph{Trust model for random sources.}
Random sources are trusted to follow the declared (or assumed
uniform) distribution.  There is no penalisation, no quit, no
slash; if a \TT{random} actor commits an out-of-support value and
refuses to reveal, the game stalls with funds held in the
contract.  A guard on a chance node is interpreted as standard
probabilistic conditioning, \TT{$\sim D$ where $\varphi$} meaning
``the value is drawn from the conditional distribution $D \mid
\varphi$''.  The MAID emitter projects self-only guards
statically into a renormalised CPD; the Gambit emitter
conditions per trace; both reject \TT{where false}.

\paragraph{Independent public moves.}  Multiple roles can appear in a
single statement.  Strategically, these are simultaneous moves:
no role is allowed to condition its choice on another role's
choice from the same statement.
\begin{lstlisting}[style=vg]
yield or split A(c: bool) B(c: bool);
\end{lstlisting}
This is exactly the situation that would invite front-running if
written naively, and is the canonical input to the automatic
commit-reveal pass described in \S\ref{sec:commit-reveal}.

\subsection{Guards and views}
\label{sec:lang:guards}

A guard is a \TT{where} clause attached to an action.  Its scope is
the snapshot of fields visible \emph{just before} the action fires,
plus the field being written.  Two distinctions matter.

\emph{Self-only versus contextual guards.}  A guard that references
only the field being written (\TT{where x != 2}) constrains the
domain of that field; a guard that references prior fields
(\TT{where Host.goat != Guest.d}) constrains the legality of a move
based on the public history.  Both are allowed.  Both are
statically required to reference only fields the actor can observe
at this point.

\emph{Public versus private fields.}  An actor cannot observe
another actor's hidden field; therefore a guard cannot read it.
The compiler rejects programs that try (it cannot enforce a
constraint nobody can check).  When the guard is on a
\TT{commit}, however, the actor's own freshly written field is
private to them: the constraint is recorded and re-checked when
the value is revealed, so that any later observer can verify the
constraint at reveal time.

\subsection{Quit handlers}

Real protocols have to cope with a player who stops responding.
Vegas treats this as a strategic move named \TT{Quit}: at any point
where a role is expected to act, the role can choose to abandon the
game.  What happens next is specified by an \emph{or-handler} on
the action:
\begin{lstlisting}[style=vg]
yield or split Guest(d: door);
yield or burn  A(c: bool);
yield or null  Odd(c: bool) Even(c: bool);
\end{lstlisting}
The handlers do the following:
\begin{itemize}\itemsep1pt
\item \TT{split} --- the quitter's deposit is split among the
  surviving roles.  In a two-player game, this is forfeit.
\item \TT{burn} --- the quitter's deposit is destroyed (or sent to
  a designated burn address).
\item \TT{null} --- the missing value becomes \TT{null} (typed
  \TT{opt T}); downstream actions and the \TT{withdraw} clause are
  required to handle the optional case.
\end{itemize}
Operationally the handler fires when a timeout elapses without the
expected action.  At the game-theoretic level, \TT{Quit} is just
another move available to the actor; equilibrium analyses see it
explicitly.  \S\ref{sec:liveness} treats the subgame structure in
detail.

\subsection{Withdraw and conservation}

A game terminates by reaching a \TT{withdraw} clause, which
allocates the entire locked pool.  Outcomes are
brace-delimited blocks of \TT{role -> exp} entries; an outcome
may additionally carry a single \TT{burn N} item that records
how much of the pool is destroyed (or, in the EVM emission,
retained in the contract permanently).  \TT{withdraw} can be
conditional:
\begin{lstlisting}[style=vg]
withdraw (A.c && B.c) ? { A -> 100; B -> 100 }
       : (A.c)        ? { A -> 200; B -> 0   }
       : (B.c)        ? { A -> 0;   B -> 200 }
       :                { A -> 90;  B -> 110 }
\end{lstlisting}
Conservation --- that, on every reachable terminal branch, the
sum of role payouts and the burn equals the sum of strategic
deposits --- is enforced statically.  The compiler runs the same
extensive-form enumeration the Gambit backend uses, and at each
terminal checks the per-branch equation; any underpayment (funds
stuck) or overpayment (contract reverts) is reported as a static
error.  Chance roles and anonymous samples have no claim on the
pool and so do not appear in withdraw payouts: their share, if
any, is recorded as a \TT{burn}.

\subsection{Types}

Vegas types are deliberately small.  The base types are:

\begin{itemize}\itemsep1pt
\item \TT{bool};
\item finite enumerations (\TT{type door = \{0, 1, 2\}}) and
  finite ranges (\TT{type bid = \{0..10\}});
\item \TT{address} (used only in joins and metadata; not first-class
  in payoff expressions);
\item \TT{int} (mapped to a 256-bit integer in the EVM backends; the
  Gambit backend rejects unbounded \TT{int} in moves because it
  needs a finite move alphabet).
\end{itemize}

Two type modifiers --- \TT{hidden} and \TT{opt} --- are
introduced by the compiler rather than written in source.  A
\TT{commit} action's parameter is implicitly typed
\TT{hidden T}: the value is private to the committing actor
until the matching \TT{reveal} fires.  A field associated with
an \TT{or null} handler is implicitly typed \TT{opt T}, marking
that it may be \TT{null} when the actor takes the \TT{Quit}
branch; the language requires \TT{withdraw} expressions to
handle the \TT{null} case explicitly.

\subsection{Macros}

Macros are hygienic expression-level abbreviations, inlined before
the IR is built.  They simplify reasoning about non-trivial
mechanisms; Vickrey auctions, for instance, are easier to read
with a \TT{secondMax} macro.

\begin{lstlisting}[style=vg]
macro max2(a: bid, b: bid): bid = (a >= b ? a : b);
macro secondMax(a: bid, b: bid, c: bid): bid =
    (a == max3(a,b,c)) ? max2(b, c)
  : (b == max3(a,b,c)) ? max2(a, c)
  :                      max2(a, b);
\end{lstlisting}

Macros expand syntactically; recursion is disallowed.  Macro
expansion is purely textual at the expression level: macros do
not introduce new actions or rewrite the action sequence.

\subsection{A walked example: Vickrey auction}
\label{sec:lang:vickrey}

Putting the constructs together, here is a three-bidder sealed-bid
second-price (Vickrey) auction with a fixed-domain bid type.  The
\TT{yield or split} statement carries three actor heads, declaring
the three bids as independent public choices.

\begin{lstlisting}[style=vg]
type bid = {0 .. 2}

macro max2(a: bid, b: bid): bid = (a >= b ? a : b);
macro max3(a: bid, b: bid, c: bid): bid = max2(max2(a, b), c);
macro secondMax(a: bid, b: bid, c: bid): bid =
    (a == max3(a,b,c)) ? max2(b, c)
  : (b == max3(a,b,c)) ? max2(a, c)
  :                      max2(a, b);
macro isWinner(myBid: bid, b1: bid, b2: bid, b3: bid): bool =
    (myBid == max3(b1, b2, b3));

game main() {
  join Seller() $ 100 B1() $ 100 B2() $ 100 B3() $ 100;

  yield or split
    B1(b: bid)
    B2(b: bid)
    B3(b: bid);

  withdraw {
    Seller -> 100 + secondMax(B1.b, B2.b, B3.b);
    B1 -> isWinner(B1.b, B1.b, B2.b, B3.b)
            ? (100 - secondMax(B1.b, B2.b, B3.b)) : 100;
    B2 -> isWinner(B2.b, B1.b, B2.b, B3.b)
            ? (100 - secondMax(B1.b, B2.b, B3.b)) : 100;
    B3 -> isWinner(B3.b, B1.b, B2.b, B3.b)
            ? (100 - secondMax(B1.b, B2.b, B3.b)) : 100;
  }
}
\end{lstlisting}

The compiler does not assume the front-running problem has been
solved.  It reads three independent public moves and rewrites the
EventGraph so that each bidder commits to a hash first and reveals
later, with every reveal gated on every other bidder's commit.
The resulting Solidity contract implements a sealed-bid auction;
the resulting Gambit game has each bidder choose a hidden value in
a common information set.  Both come from the same source.

\subsection{Grammar}

For reference, the full surface grammar lives in the ANTLR file
\file{Vegas.g4} at the project root; the salient productions are
roughly:
\begin{align*}
\mathit{program} &::= (\mathit{typeDec} \mid \mathit{macroDec})^*\ \TT{game}\ \mathit{ident}\ \TT{()}\ \TT{\{}\ \mathit{ext}\ \TT{\}}\\
\mathit{ext} &::= \mathit{kind}\ (\texttt{or}\ \mathit{handler})?\ \mathit{query}^+\ \TT{;}\ \mathit{ext}\\
             &\mid \TT{random}\ \mathit{role}\ (\TT{()})?\ \TT{;}\ \mathit{ext}\\
             &\mid \TT{sample}\ \TT{(}\ \mathit{varDec}^+\ \TT{)}\ \TT{;}\ \mathit{ext}\\
             &\mid \kwithdraw\ \mathit{outcome}\\
\mathit{kind} &::= \kjoin \mid \kyield \mid \kreveal \mid \kcommit\\
\mathit{query} &::= \mathit{role}\ \TT{(}\ \mathit{paramDec}^*\ \TT{)}\ (\$\,\mathit{int})?\ (\kwhere\ \mathit{exp})?\ (\TT{||}\ \mathit{queryHandler})?\\
\mathit{paramDec} &::= \mathit{name}\ \TT{:}\ \mathit{type}\ (\TT{$\sim$}\ \mathit{dist})?\\
\mathit{dist} &::= \TT{uniform \{ \dots \}} \mid \TT{weighted \{ }\mathit{v}\TT{:}\,\mathit{w},\dots\ \TT{\}}\\
\mathit{handler} &::= \TT{split} \mid \TT{burn} \mid \TT{null}\\
\mathit{item} &::= \mathit{role}\ \TT{->}\ \mathit{exp} \mid \TT{burn}\ \mathit{exp}
\end{align*}
The \TT{random} production does not admit \TT{\$ deposit}: chance
actors do not stake.  The grammar is unsurprising; what makes
Vegas work is the analysis done in the next several sections.
